\documentclass[addpoints]{exam}
% leqno 將方程式編號放在左側

%\printanswers

\usepackage[top=2.5cm,bottom=2cm,left=2.5cm,right=2.5cm,headsep=10pt,a4paper]{geometry} % 引用頁面幾何套件，設定頁面邊界

\usepackage{siunitx} % 角度
% \input{NewCommands}

\usepackage{amsmath,amssymb,amsthm} % 引用常見的數學符號與設定


\theoremstyle{definition}
\newtheorem*{definition}{Definition}
\newtheorem*{theorem}{Theorem}
\newtheorem*{remark}{Remark}
\newcommand{\R}{\mathbb{R}}

\usepackage{lastpage} % 取得最後頁碼




\usepackage{graphicx} % 引用圖檔內嵌入套件
\graphicspath{{Pictures/}} % 設定圖檔目錄
\usepackage{xcolor} %引用顏色套件
\usepackage{enumerate}
\usepackage{float}
\usepackage{hyperref}
\usepackage{mdwlist} % use suspend & resume to successive enumerate

\usepackage{CJKutf8} % 設定中文

\def \lflr{\left\lfloor} %自定義 floor function
\def \rflr{\right\rfloor} %自定義

% question separation
\renewcommand{\questionshook}{%
  \setlength{\itemsep}{0.5cm}%
}

\allowdisplaybreaks %equation allowed to next page

\firstpageheader{\today}
                {}
                {}
\footer{}{\sf Page \thepage~of~\pageref{LastPage}}{}

\begin{document}

\begin{CJK*}{UTF8}{gbsn}
% Title
\begin{center}
    {\Large{\bf Introduction to Mathematical Analysis II\\
    Homework 5  Due  April 4 (Saturday), 2026\\
    Please submit your group homework online in PDF format.\\
You are expected to work collaboratively within your group. 
Please do not obtain complete solutions directly from AI tools. 
The purpose of this assignment is to develop your own mathematical reasoning and problem-solving skills. %  Homework X Brief Solution
  }} \\ 
    \large
    %\text{Due date: 3/9/2023}
\end{center}

\noindent\rule{16.2cm}{0.4pt}


\begin{questions}


% ============ Question x ============
\question  (20 pts)  (Exercises 7.2.2.)
Let $A$ be a subset of $\mathbf{R}^n$, and let $B$ be a subset of $\mathbf{R}^m$.
Note that the Cartesian product $\{(a,b): a\in A,\, b\in B\}$ is then a subset of
$\mathbf{R}^{n+m}$. Show that
$$
m^*_{n+m}(A\times B)\le m^*_n(A)\,m^*_m(B).
$$
Here we adopt the convention that $c\times +\infty=+\infty\times c$ is equal to $+\infty$
for any $0<c\le +\infty$ and equal to zero for $c=0$.
(It is in fact true that $m^*_{n+m}(A\times B)=m^*_n(A)m^*_m(B)$, but this is
substantially harder to prove.)

\begin{solution}
\end{solution}

\question (15 pts) (Exercises 7.2.3.)


\begin{itemize}
\item[(a)] Show that if $A_1\subseteq A_2\subseteq A_3\subseteq\cdots$ is an increasing sequence
of measurable sets (so $A_j\subseteq A_{j+1}$ for every positive integer $j$), then we have
$$
m\Bigl(\bigcup_{j=1}^\infty A_j\Bigr)=\lim_{j\to\infty} m(A_j).
$$
\item[(b)] Show that if $A_1\supseteq A_2\supseteq A_3\supseteq\cdots$ is a decreasing sequence
of measurable sets (so $A_j\supseteq A_{j+1}$ for every positive integer $j$), and $m(A_1)<+\infty$,
then we have
$$
m\Bigl(\bigcap_{j=1}^\infty A_j\Bigr)=\lim_{j\to\infty} m(A_j).
$$
\end{itemize}

\begin{solution}

\begin{enumerate}
\item[(a)] 



\bigskip
\item[(b)] 




\end{enumerate} 
\end{solution}



\question (15 pts)
Give examples showing that:
\begin{enumerate}
    \item a set can have boundary of outer measure $0$;
    \item a set can have boundary of positive outer measure.
\end{enumerate}

\begin{solution} 

\end{solution}

\question (20 pts) 
Recall that a function $f: [a, b] \to \mathbf{R}$ is \textit{Lipschitz} if there exists $L \geq 0$ such that $|f(x) - f(y)| \leq L|x-y|$ for all $x, y \in [a, b]$. 
Define the graph of $f$ as:
\[
\Gamma_f := \{(x, f(x)) : x \in [a, b]\} \subset \mathbf{R}^2.
\]
Prove that $m^*_2(\Gamma_f) = 0$. 
\textit{Hint: Partition $[a, b]$ into $n$ sub-intervals and cover the corresponding pieces of the graph with rectangles. Consider the limit as $n \to \infty$.}

\begin{solution} 
 
 \end{solution}
 
\question (15 pts) Let $A\subset \mathbf{R}$. Show that
\[
m_2^*(A\times\{0\})=0
\]
in $\mathbf{R}^2$.
More generally, if $m_1^*(A)=0$, prove that
\[
m_2^*(A\times[0,1])=0.
\]

\medskip
\noindent\textbf{Hint.}
Start from interval coverings of $A$, then build rectangles in $\mathbf{R}^2$.
\begin{solution} 

\end{solution}

\question (15 pts) Let $\{P_j\}_{j=1}^\infty$ be a sequence of properties defined for points $x\in \mathbf{R}^n$. For each $j\in \mathbf{N}$, define the exceptional set
\[
E_j:=\{x\in \mathbf{R}^n : P_j(x)\text{ does not hold}\}.
\]
We say that $P_j$ holds \emph{outer-almost everywhere} if
\[
m^*(E_j)=0.
\]

Assume that each property $P_j$ holds outer-almost everywhere. Prove that the property
\[
P(x): \quad P_j(x)\text{ holds for every } j\in \mathbf{N}
\]
also holds outer-almost everywhere.

Equivalently, show that the set
\[
E:=\{x\in \mathbf{R}^n : P_j(x)\text{ holds for every }j\in \mathbf{N}\}
=\bigcap_{j=1}^\infty \{x\in \mathbf{R}^n : P_j(x)\text{ holds}\}
\]
has complement of outer measure zero; that is,
\[
m^*(\mathbf{R}^n\setminus E)=0.
\]


\begin{solution} 

\end{solution}

\end{questions}
\end{CJK*}
\end{document}

